\section{Schemes, weights, and the limits of the comparison}
\label{sec:framework}

All arithmetic cohomology over the rational integers in this paper lives on
$U_\ell=\operatorname{Spec}\mathbb Z[1/\ell]$. The Tate lattice
$\mathbb Z_\ell(r)$ is the inverse system of finite etale sheaves
$\mu_{\ell^n}^{\otimes r}$, interpreted by continuous cohomology.
The discrete sheaf $\mathbb Q_\ell/\mathbb Z_\ell(-e)$ is its corresponding
divisible torsion coefficient object. These coefficient conventions,
the local conditions, and the compact/discrete duality are part of the
statements; replacing them by rational cohomology loses the finite groups
being bounded.

The weight convention is Deligne's \emph{Weil II}, Definitions 1.2.1--1.2.2,
p.~153, and 1.2.5(iv), p.~154 \cite{DeligneWeilII}: $\mathbb Q_\ell(1)$ has weight $-2$.
Thus the rational Tate object $\mathbb Q_\ell(-e)$ has weight $2e$,
and $\mathbb Q_\ell(e+1)$ has weight $-2(e+1)$. These statements refer
to rational sheaves and geometric Frobenius. They do not assign a
Frobenius purity theorem to a torsion module or prove its vanishing.
The finite $K$-groups at issue have zero rationalization even when their
integral vanishing is unknown. This elementary observation explains why
a weight or sign argument requires the additional integral congruence
and duality statements reconstructed above.

For the archive passage, \emph{SGA 4}, IX.3.2--3.3 supplies the
prime-invertible Kummer sequence and Hilbert 90. The first Chern class
is the Kummer boundary of a line bundle; \emph{SGA 5}, VII,
Proposition 3.4(i), identifies the degree-one term of the total Chern
class with that boundary \cite{SGA5}. The finite root-coordinate gluing proof in
the archive section fixes the precise form needed here. The tame
cover comparison is \emph{SGA 1}, XIII.5.2--5.6, with X.3.6 for the
discrete-valuation Abhyankar lemma. The extension of regular functions
across codimension at least two uses \emph{EGA} $\mathrm{IV}_2$,
5.8.6, 5.9.8--5.9.9 and 5.10.5. Recognition with the full generic
arithmetic marking additionally uses Mochizuki's anabelian theorems;
it is not a consequence of an FGA representability statement.

The $\lambda$-ring language belongs first to $K_0$: exterior powers
define $\lambda^i$, and the determinant or first gamma-filtration
piece records line bundles. An exact locator is \emph{SGA 6}, V,
Definition 2.1, Example 2.2(b), and Lemma 2.2.1 \cite{SGA6}. Higher $K_{4k}$-groups
require the higher etale comparison used in the $K$-group section. Neither
the $K_0$ operations nor a first Chern class is that comparison map.

Crystals and divided powers concern a different coefficient setting,
especially reduction in residue characteristic $\ell$. No new theorem
about a crystal, a divided-power envelope, or a crystalline comparison
is proved in this run. The etale Kummer calculation deliberately removes
$\ell$ from the base; for noninvertible exponents one must use flat
Kummer theory and then state a separate comparison with any crystalline
or prismatic object. The inseparable counterexample in the archive
section proves why the etale assertion itself cannot be transported
unchanged. Assigning one of the present finite Selmer estimates to a
crystalline object without such a map would add an unproved assertion.

Finally, the rational mixed-Tate category over $\operatorname{Spec}
\mathbb Z$ is the category used by Deligne--Goncharov and Brown. Its
base scheme controls ramification; its rational coefficients control
which extensions remain visible. The integral arithmetic lattice of a
pro-$\ell$ image is an additional object. The counterexample at $37$
shows that these distinctions change the answer to B3.

\subsection{Location of the new deductions in these frameworks}

The following grouping specifies the geometric object and the exact
framework used by each family of new findings. A source in the last
column supplies the stated language or comparison; it is not credited
with a later arithmetic or computational conclusion that it does not
contain.

\small
\begin{longtable}{>{\raggedright\arraybackslash}p{.23\textwidth}>{\raggedright\arraybackslash}p{.34\textwidth}>{\raggedright\arraybackslash}p{.34\textwidth}}
\toprule
Finding family & Geometric object & Precise framework and scope\\
\midrule
\endhead
K17-S17, HOST, PT-Q; VK-EXP7 and POLY-CERT &
Strict-local-condition cohomology of
$\operatorname{Spec}\mathcal O_E[1/\ell]$ with negative Tate
coefficients; higher $K$-theory of $\operatorname{Spec}\mathbb Z$. &
SGA~4 IX.3.2--3.3 fixes finite Kummer coefficients; continuous limits
and arithmetic duality are the separately cited Jannsen/Milne inputs.
Weil~II 1.2.1--1.2.5 fixes rational weights, not torsion vanishing.\\
K17-S17-UNIT, U3 and GL2 models; K12 and higher vanishings; REFLECTOR and Bernoulli certificates &
Integral comparison of Frobenius traces in finite-dimensional algebraic modular forms, modular-curve representations, and strict torsion cohomology of $\operatorname{Spec}\mathbb Z[1/\ell]$. &
SGA~4 IX.3.2--3.3 places the finite Kummer radicals and Tate coefficients; Weil~II 1.2.1--1.2.5 fixes weights. The local monodromy and crystalline comparison are specifically Saito \S2c Theorem~3 and Colmez--Fontaine \S3.3, Propositions~4.2,4.4, as cited above. No SGA theorem alone supplies the finite congruence calculation.\\
K17-HOST-GREENBERG-GATE &
Norm towers of finite arithmetic cohomology and finitely generated
Iwasawa modules. &
The same finite etale coefficient category; the countermodels test
growth inequalities only. No crystal or divided-power object is
identified with these modules.\\
K17-P2-MOMENT, VANDIVER-KERNEL, GREENBERG-KERNEL;
CUP-EXT and RATIONAL-MAGNUS &
The arithmetic action on the geometric pro-$\ell$ fundamental group
of $\mathbb P^1-\{0,1,\infty\}$; its weight filtration and
finite-level cup products. &
SGA~1 XIII.5.2--5.6 controls tame cusp inertia. SGA~4 IX.3.2--3.3
and the stated low-degree etale/group comparison place the cups.
The mixed-Tate rational target is Deligne--Goncharov \S1.6,
Propositions~1.7,1.9,2.3 and \S2.4, with Brown's cited theorem.\\
K17-OUT-SAT, 2-ALL and the two integral classifications &
An actual $\mathbb Z_\ell$ Lie lattice inside the rational
Tannakian image, with finite torsion in its abelianization. &
The above mixed-Tate framework supplies the rational comparison.
The integral lattice and its even torsion are additional arithmetic
data; SGA~6 V.2.1--2.2.1 on $K_0$ does not identify them.\\
K17-K12-THREE-RANKS, DOUBLE, CIRCUIT-SIMPLE,
ORIENT-SURVIVOR and SURVIVOR-CYCLE &
Quillen's rank filtration for $K(\operatorname{Spec}\mathbb Z)$,
Steinberg coefficients and integral Voronoi chains. &
These use the exact Quillen, Voronoi, building and Hurewicz sources
cited in their proofs. SGA~6 V describes the degree-zero
$\lambda$-ring; no claim identifies its operations with this
higher-homology computation.\\
F149-KUM, FULLINJ, CHERN and RELATIVE &
The line at a marked boundary section, its nonzero-vector scheme,
and the relative cusp-normalizer extension. &
Cote~149 scans~79--82 and~85; SGA~4 IX.3.2--3.3;
SGA~5 VII.3.4(i); SGA~1 XIII.4.1,4.4 and Deligne~15.19,15.26.
Normal and conormal orientations are kept distinct.\\
F149-GENERIC, MARKED-RECOGNITION and BARE-EXTENSION &
Geometric tangent splittings inside the full continuous
$H^1$ torsor, with the generic Galois marking retained. &
Cote~149 scans~83--85; EGA~$\mathrm{IV}_2$~5.10.5 for extension
from codimension one; finite Kummer at every exponent;
Mochizuki Theorems~A and~B for the extra anabelian reconstruction.
The isolated extension fails by the explicit cocycle automorphism.\\
\bottomrule
\end{longtable}
\normalsize

